The branching ratio $\branchingRatio$ is the number one quotes to say how self-excited a
process is.
This section says what it, and the quantities that travel with it, actually tell us.
Throughout, $\excitation \geq 0$, $\baseIntensity \geq 0$ and $\baseIntensity \neq 0$,
and $\branchingMatrix = \excitation/\decay$ as in \eqref{eq.GammaIsAOverBeta}.

We shall use three facts about non-negative matrices,
all of them consequences of the Perron--Frobenius theorem.
For $\branchingMatrix \geq 0$:
the branching ratio $\branchingRatio$ is itself an eigenvalue of $\branchingMatrix$, and it is real;
it lies between the smallest and the largest column sum,
so constant column sums $c$ force $\branchingRatio = c$;
and if $\branchingRatio < 1$ then
$(I-\branchingMatrix)^{-1} = \sum_{k\geq 0} \branchingMatrix^{k}$
converges with non-negative entries.

\subsection{What the branching matrix says}

By Proposition \ref{prop.clusterRepresentation},
$\branchingMatrix\subscriptee$ is the expected number of direct offspring of type $e$
born of one event of type $\eone$.
Its column sums are therefore readable one by one,
without any spectral computation,
and they are the most informative single summary of a fitted or specified kernel.
For the six-type order flow used in these notes they are
\begin{center}
 \begin{tabular}{lcccccc}
  \toprule
  $\eone$ & market & market & limit & limit & withdraw & withdraw \\
  & buy & sell & buy & sell & bid & ask \\
  \midrule
  $\sum_e \branchingMatrix\subscriptee$ & $1.425$ & $1.425$ & $0.199$ & $0.199$ & $1.036$ & $1.036$ \\
  \bottomrule
 \end{tabular}
\end{center}
One market order begets $1.43$ direct events of one type or another,
a withdrawal $1.04$, and a limit order only $0.20$.
That single line is the replenishment mechanism of the model expressed as a number:
the two kinds of event that deplete a queue --- an aggression against it and a withdrawal
from it --- both beget successors, and what they chiefly beget is the limit orders that
refill the side they emptied, while a limit order that has already refilled begets little.
Note again, per Section \ref{sec.exponentialKernel},
that the corresponding column sums of $\excitation$ are $\decay$ times these,
and are jumps in intensity rather than counts.

\subsection{Stationarity and the mean intensity}

\begin{prop}\label{prop.stationarity}
 Let $\hawkesKernel \geq 0$ be locally integrable with $\branchingMatrix$ finite,
 and let $\baseIntensity \geq 0$, $\baseIntensity \neq 0$.
 A stationary version of the Hawkes process of Definition \ref{def.hawkesprocess}
 with finite mean intensity exists if and only if $\branchingRatio < 1$,
 and in that case
 \begin{equation}\label{eq.stationaryIntensity}
  \stationaryIntensity := \Expectation \big[ \intensity[ ](t) \big]
  = (I - \branchingMatrix)^{-1} \baseIntensity.
 \end{equation}
\end{prop}

Taking expectations in \eqref{eq.intensity_ordHawkes} at stationarity gives
$\stationaryIntensity = \baseIntensity + \branchingMatrix \stationaryIntensity$,
which is \eqref{eq.stationaryIntensity};
what $\branchingRatio<1$ supplies is that the Neumann series converges,
so that the solution exists and is non-negative.
We write
\begin{equation*}
 \totalRate := \mathbf{1}^{\transpose} \stationaryIntensity,
 \qquad
 \bar{\baseIntensity} := \mathbf{1}^{\transpose} \baseIntensity
\end{equation*}
for the total event rate and the total immigration rate.

Read in the other direction, \eqref{eq.stationaryIntensity} is how one specifies a model:
choose the total rate one wants to observe, choose the shape of the excitation and the
branching ratio, and let $\baseIntensity = (I - \branchingMatrix)\stationaryIntensity$ follow.
The baselines are then an output, and the constraint $\baseIntensity \geq 0$ is a genuine
restriction on the pairs $(\stationaryIntensity, \branchingMatrix)$ one may ask for.

\subsection{Clusters, and how much of the flow is endogenous}

The events partition into clusters, one per immigrant.
Write $C_j$ for the expected size of the cluster founded by an immigrant of type $j$.
Summing the generations of the Galton--Watson process of
Proposition \ref{prop.clusterRepresentation},
\begin{equation}\label{eq.clusterSize}
 C_j = \mathbf{1}^{\transpose} (I - \branchingMatrix)^{-1} e_j .
\end{equation}
The $k$-th term of the series counts the $k$-th generation,
and the term $k=0$ is the immigrant itself:
$C_j$ counts the ancestor.
The expected number of its strict descendants is $C_j - 1$.
Averaging over the type of the immigrant --- weighted by $\baseIntensity_j / \bar{\baseIntensity}$,
since it is immigrants that found clusters --- gives
\begin{equation}\label{eq.meanClusterSize}
 \bar{C} = \sum_j \frac{\baseIntensity_j}{\bar{\baseIntensity}} C_j
 = \frac{\mathbf{1}^{\transpose}(I-\branchingMatrix)^{-1}\baseIntensity}{\bar{\baseIntensity}}
 = \frac{\totalRate}{\bar{\baseIntensity}} .
\end{equation}
The identity $\bar{\baseIntensity} \, \bar{C} = \totalRate$ is worth stating on its own:
immigrants arrive at rate $\bar{\baseIntensity}$, each brings $\bar C$ events,
and every event belongs to exactly one cluster.
It is the cheapest available check that an implementation of
\eqref{eq.stationaryIntensity} and \eqref{eq.clusterSize} is consistent.

The complementary quantity is the \emph{endogenous fraction}
\begin{equation}\label{eq.endogenousFraction}
 1 - \frac{\bar{\baseIntensity}}{\totalRate},
\end{equation}
the long-run share of events that are offspring rather than immigrants.
It is a ratio of rates, not a probability attached to any single event,
and by \eqref{eq.meanClusterSize} it equals $1 - 1/\bar C$.

\begin{remark}[what $\branchingRatio$ is, and what it is not]\label{remark.rhoIsNot}
 It is tempting to read $\branchingRatio$ as the endogenous fraction,
 and $1/(1-\branchingRatio)$ as the mean cluster size.
 The two readings are equivalent to each other, by \eqref{eq.meanClusterSize},
 and they hold exactly when
 $\totalRate = \bar{\baseIntensity}/(1-\branchingRatio)$.
 That is the case in particular when
 $\mathbf{1}^{\transpose}\branchingMatrix = \branchingRatio \mathbf{1}^{\transpose}$,
 the column sums of $\branchingMatrix$ being constant,
 and in particular when $\branchingMatrix\baseIntensity = \branchingRatio\baseIntensity$,
 the baselines being proportional to a right Perron vector.
 Neither is generic, and neither holds for the parametrisation of these notes:
 there the column sums are $(1.43, 1.43, 0.20, 0.20, 1.04, 1.04)$,
 and at $\branchingRatio = 0.60$ the endogenous fraction is $0.682$ and
 $\bar C = 3.149$, against $1/(1-\branchingRatio) = 2.5$.
 The gap is small here and need not be;
 the point is that $\branchingRatio$ is a spectral radius,
 while \eqref{eq.meanClusterSize} and \eqref{eq.endogenousFraction} are averages against
 $\baseIntensity$, and a spectral radius does not know about $\baseIntensity$.
 Only the stronger statement --- that $C_j = 1/(1-\branchingRatio)$ for \emph{every} $j$ ---
 genuinely requires constant column sums.
\end{remark}

\subsection{Timescales}

The mean response of the state to a single event is
$\mathbf{1}^{\transpose} e^{(\excitation - \decay I)t} e_j$,
a sum of $\numEventTypes$ exponentials at the rates $\decay(1-\gamma_i)$,
$\gamma_i$ running over the eigenvalues of $\branchingMatrix$.
Since $\branchingMatrix \geq 0$, the branching ratio is a real eigenvalue and
$\decay(1-\branchingRatio)$ is the smallest of these rates.
We therefore call
\begin{equation}\label{eq.relaxationTime}
 \frac{1}{\decay(1-\branchingRatio)}
\end{equation}
the relaxation time of the process.
It is the time constant of the \emph{slowest} mode,
which is to say that the response is asymptotically proportional to
$e^{-\decay(1-\branchingRatio)t}$;
it does not describe the response at short times,
and it is not a half-life.
On the parametrisation of these notes the response to a market order rises above one before
it falls.

Two dimensionless groups organise everything that follows.
The first, $\totalRate/\decay$, is the number of events that arrive within one excitation
timescale: it says whether clustering is visible at the rate the process is sampled.
The second, $\totalRate w$, is the number of events in an observation window of length $w$,
and it says whether a window sum is a statistic or a single event.
A study conducted at $\totalRate/\decay \ll 1$ measures a process whose self-excitation has
expired between consecutive events, and finds, correctly, that there is nothing there.

\subsection{Sub- and supercriticality}

Everything above assumes $\branchingRatio < 1$.
It is worth being precise about what fails otherwise, because it is less than one expects.

For a linear Hawkes process and every finite horizon $\timeHorizon$,
$\Expectation[\countingProc(\timeHorizon)^2] < \infty$,
so on $[0,\timeHorizon]$ the process
$\martingale := \multiCountingProc - \compensator$
is a square-integrable martingale,
with $\langle \martingale_e \rangle = \compensator_e$
and $\langle \martingale_e, \martingale_{\eone} \rangle \equiv 0$ for $e \neq \eone$,
the latter by \eqref{eq.noCommonJumps}.
This holds at every $\branchingRatio$, sub- or supercritical.
Subcriticality is not what makes $\multiCountingProc - \compensator$ a martingale.

What subcriticality buys is three other things:
a stationary version exists, by Proposition \ref{prop.stationarity};
the second moments of the \emph{state} $(\decayedCounts, \intensity[ ])$ are bounded uniformly
in $\timeHorizon$, which is the subject of Section \ref{sec.secondOrder};
and the process is ergodic, so that time averages along one path estimate the stationary
quantities above.
Each of these fails at $\branchingRatio \geq 1$, where the intensity grows without bound and
the event count grows exponentially.
Note that the martingale is not bounded in $L^2$ on $[0,\infty)$ even when
$\branchingRatio<1$:
$\Expectation[\martingale_e(\timeHorizon)^2] = \Expectation[\compensator_e(\timeHorizon)]
\sim \stationaryIntensity[e] \timeHorizon$ diverges.
The boundedness that subcriticality delivers is boundedness of the state, not of the martingale.
